\chapter{Mixture of Experts——参数效率的极致}
\label{ch:moe}
%% ============================================================

\section{核心思想：不是每个神经元都需要工作}

想象一个公司有 256 名专家，分别精通不同领域。
当一个任务到来时，你不需要让所有 256 人都参与——
你只需要挑选最相关的 8 个人来处理。
这就是 Mixture of Experts (MoE) 的核心思想。

\subsection{历史脉络}

MoE 的思想远比深度学习更古老。
1991 年，Jacobs 等人提出了最早的 MoE 框架——% NEW REF: jacobs1991moe = Adaptive Mixtures of Local Experts, Jacobs et al., Neural Computation 1991
用多个「专家」网络分别处理输入空间的不同区域，
一个「门控」网络决定每个输入由哪些专家处理。
但这个想法在之后的 25 年里基本停留在学术论文中，
没有被大规模应用。

2017 年，Shazeer 等人~\cite{shazeer2017moe}
将 MoE 引入了 NLP 领域，
在 LSTM 语言模型中嵌入了稀疏门控的 MoE 层，
用 $>$1000 个专家实现了当时的最大模型（137B 参数）。
这篇论文的标题直接叫做
``Outrageously Large Neural Networks''——
「大得离谱的神经网络」。

2022 年，Switch Transformer~\cite{fedus2022switch}
将 MoE 扩展到 Transformer 架构，
用 top-1 路由（每个 token 只发给 1 个专家）
训练了一个 1.6T 参数的模型。
Switch Transformer 证明了一个关键点：
\textbf{即使只激活一个专家，
只要专家数量足够多、路由足够准确，
模型性能就可以超越同规模的稠密模型}。

自此，MoE 从一种学术好奇心
变成了 LLM 训练的主流架构选择。

\subsection{稠密模型 vs 稀疏模型}

传统的「稠密」模型（Dense Model）在处理每个 token 时
使用所有参数——参数量直接决定了计算量。
MoE 模型则打破了这种绑定：
\textbf{总参数量决定了模型的知识容量，
但每个 token 的计算量只取决于被激活的参数}。

DeepSeek-V3 的数据最能说明这种效率：
总参数 671B（知识容量巨大），
每个 token 仅激活 37B（计算成本低），
激活率约 5.5\%。

用一个类比来理解：
稠密模型像一个通才——不管什么问题，
他调动全部知识和精力来回答，效果不错但很累。
MoE 模型像一个高效的团队——
不管什么问题，先判断这是什么类型的问题，
然后把最擅长的几个人召集起来处理，
其他人可以休息。
团队的总知识量（总参数）可以很大，
但每次实际消耗的精力（计算量）很小。

\section{路由机制}

MoE 层替换了 Transformer 中的 FFN 层。
对于输入 token $\mathbf{x}$，
路由器（一个简单的线性层）计算每个专家的分数，
然后选择得分最高的 $K$ 个专家：
\begin{equation}
  g_i = \frac{e^{s_i}}{\sum_{j \in \mathrm{TopK}} e^{s_j}}, \quad
  s = W_{\mathrm{router}} \mathbf{x}, \quad
  \mathbf{y} = \sum_{i \in \mathrm{TopK}} g_i \cdot \mathrm{Expert}_i(\mathbf{x})
  \label{eq:moe-routing}
\end{equation}

路由器本身极其轻量——
它只是一个 $[d, N_{\mathrm{experts}}]$ 的线性投影，
没有偏置，没有激活函数。
它的参数量微不足道（$4096 \times 256 \approx 1M$），
但它做出的决策至关重要：
将每个 token 送到最合适的专家去处理。

训练过程中，路由器通过梯度下降与专家一起优化。
一个有趣的现象是路由器会自发学到有意义的分工：
某些专家可能专注于代码相关的 token，
某些专注于数学符号，
某些专注于特定语言的文本。
这种\textbf{自动专业化}不是人为设定的——
它完全是训练过程中涌现出来的。

\subsection{Top-K 路由}

Top-K 路由是目前最主流的方案。
$K$ 的选择涉及质量与效率的权衡：

\textbf{Top-1}（Switch Transformer~\cite{fedus2022switch}）：
每个 token 只发给 1 个专家。
计算效率最高，但单个专家可能无法充分处理复杂 token。
而且 top-1 的路由决策更「尖锐」——
一旦选错专家，没有其他专家来补救。

\textbf{Top-2}（Mixtral~\cite{jiang2024mixtral}、GShard）：
每个 token 发给 2 个专家，% NEW REF: lepikhin2021gshard = GShard: Scaling Giant Models with Conditional Computation, Lepikhin et al., ICLR 2021
结果按路由权重加权求和。
这是一种「保险」策略——
即使一个专家不太合适，
另一个通常能弥补。
Top-2 是最常见的选择，
在质量和效率之间取得了良好平衡。

\textbf{Top-8}（DeepSeek-V3~\cite{deepseekv3}）：
每个 token 发给 8 个专家。
看起来计算量更大？
不——因为 DeepSeek 使用了 256 个\textbf{细粒度小专家}，
每个专家只有标准专家的 $\sim$1/32 大小。
选 8 个小专家的总计算量与选 2 个大专家相当，
但组合空间爆炸式增长（从 28 种增加到 44 亿种）。

\subsection{Token Choice vs Expert Choice}

上述路由方式都是\textbf{Token Choice}——
每个 token 独立选择自己要去的专家。
这会导致一个固有问题：
热门专家可能收到过多 token，冷门专家闲置。

Zhou 等人（2022）提出了\textbf{Expert Choice} 路由：% NEW REF: zhou2022expertchoice = Mixture-of-Experts with Expert Choice Routing, Zhou et al., NeurIPS 2022
反过来，让专家选择 token。
每个专家从所有 token 中选取得分最高的 $k$ 个来处理。
这天然保证了负载均衡——
每个专家处理相同数量的 token，不存在过载或空闲。

但 Expert Choice 有一个问题：
某些 token 可能被多个专家选中（计算冗余），
而某些 token 可能不被任何专家选中（信息丢失）。
在实践中，Token Choice + 负载均衡机制
仍然是更主流的方案。

\subsection{其他路由策略}

\textbf{Hash 路由}是最简单的基线：
根据 token 的哈希值确定性地分配到专家。
不需要可学习的路由器，负载完全均衡，
但缺乏语义感知——
代码 token 和诗歌 token 可能被分到同一个专家。
实验表明 Hash 路由的性能显著低于学习型路由。

\textbf{随机路由}也是一个有用的基线：
随机均匀地将 token 分配到专家。
在某些配置下，随机路由的表现出人意料地不差——
因为即使路由是随机的，
专家本身仍然可以通过梯度下降
学到对不同类型输入的通用处理能力。
这提示了一个有趣的问题：
路由器到底有多重要？
答案是：在大规模训练中，学习型路由\textbf{显著}优于随机路由，
但差距不如你直觉认为的那么大。

\subsection{动态路由与可微路由（2025--2026 前沿）}
\label{sec:advanced-routing}

标准的 Top-K 路由有一个根本局限：
\textbf{每个 token 固定激活 $K$ 个专家}，
无论 token 的复杂度如何。
简单的 token（如 ``the''）可能只需要 1 个专家就够了，
而复杂的 token（如数学公式中的关键符号）
可能需要 5 个以上的专家协同处理。

\textbf{DynaMoE}~(2026) 放松了这两个约束：% NEW REF: dynamoe2026 = DynaMoE: Dynamic Token-Level Expert Activation with Layer-Wise Adaptive Capacity, Guelmez, arXiv 2026
（1）每个 token 激活的专家数量根据输入复杂度
动态变化（从最少 1 个到层级相关的最大值）；
（2）不同层的专家总容量可以不同——
底层可能需要更多专家处理低级特征的多样性，
顶层可能需要更少但更专精的专家。
DynaMoE 提出了 6 种容量调度策略
（线性递增、线性递减、中间峰值等），
让模型自动学习最优的层间容量分配。

\textbf{DirMoE} 是 ICLR 2026 上提出的% NEW REF: dirmoe2026 = DirMoE: Differentiable Probabilistic Router for Mixture-of-Experts, ICLR 2026
\textbf{完全可微}的概率路由器。
传统的 Top-K 路由不可微——
``选择第 3 和第 7 个专家'' 这个离散决策
无法产生干净的梯度信号。
DirMoE 将路由分解为两个可微的概率过程：
\begin{itemize}[noitemsep]
  \item \textbf{Bernoulli 过程}决定每个专家是否被激活
        （连续松弛的 Gumbel-Sigmoid）。
  \item \textbf{Dirichlet 过程}决定被激活专家之间的权重分配。
\end{itemize}
这种分离使得路由在端到端训练中拥有干净的梯度流。
DirMoE 还引入了基于 Simpson 指数的``稀疏度旋钮''——
通过一个标量超参数精确控制期望的活跃专家数量，
\textbf{无需辅助损失}来维持负载均衡。
实验表明 DirMoE 在匹配标准 MoE 吞吐量的同时
实现了更稳定的路由和更均匀的专家利用率。

\textbf{L2R}（Low-Rank \& Lipschitz-Controlled Routing）% NEW REF: l2r2026 = L2R: Low-Rank and Lipschitz-Controlled Routing for Mixture-of-Experts, Yang et al., arXiv 2026
从另一个角度改进路由：
在低秩潜在空间中执行路由决策，
并通过 Lipschitz 控制确保路由函数的平滑性。
直觉上，标准的线性路由器在高维空间中
可能因``角度集中''现象而产生
几乎相同的分数——微小的输入变化导致路由决策翻转。
L2R 通过约束路由函数的 Lipschitz 常数，
确保了路由决策的稳定性。

用代码来看，MoE 路由器的核心逻辑可以用几行实现：

\begin{lstlisting}[caption={MoE router implementation}]
class MoERouter(nn.Module):
    def __init__(self, d_model, n_experts, top_k):
        super().__init__()
        self.gate = nn.Linear(d_model, n_experts, bias=False)
        self.top_k = top_k

    def forward(self, x):
        # x: (batch, seq_len, d_model)
        scores = self.gate(x)  # (batch, seq_len, n_experts)

        # Select top-k experts per token
        topk_scores, topk_ids = scores.topk(self.top_k, dim=-1)

        # Normalize scores (softmax over selected experts only)
        topk_weights = F.softmax(topk_scores, dim=-1)

        return topk_weights, topk_ids
\end{lstlisting}

这段代码展示了路由的三个步骤：
线性投影得到每个专家的分数，
TopK 选出得分最高的 $K$ 个专家，
在选出的专家上做 softmax 得到归一化的权重。
这些权重决定了各专家输出的贡献比例。
其中 $g_i$ 是归一化后的路由权重，
只在被选中的 TopK 专家上计算 softmax。

\section{负载均衡：MoE 训练的核心挑战}

\subsection{为什么负载不均衡是灾难性的}

MoE 训练中最棘手的问题是\textbf{专家坍缩}
（expert collapse）：
路由器学会将绝大多数 token 发送到少数几个「强」专家，
而大部分专家几乎不接收任何 token，逐渐退化为随机噪声。

为什么这是灾难性的？两个层面：

\textbf{计算浪费}：
假设有 64 个专家分布在 64 块 GPU 上，
如果 80\% 的 token 都被路由到 4 个专家，
那么只有 4 块 GPU 在满负荷工作，
另外 60 块 GPU 基本在空转——
你花了 64 块 GPU 的钱，只得到了 4 块 GPU 的效果。
在分布式训练中，整个系统的速度受最慢节点（最忙的专家）制约，
所以负载不均衡直接导致训练吞吐量暴跌。

\textbf{知识容量浪费}：
MoE 的核心优势是「用总参数量换知识容量」。
如果大部分专家闲置，等价于这些参数不存在——
671B 参数的 MoE 模型可能退化为
只有几十 B 参数的效果。

\subsection{传统方案：辅助损失}

最早的解决方案~\cite{shazeer2017moe}
是在训练目标中加入一个辅助损失来惩罚不均衡：
\begin{equation}
  L_{\mathrm{total}} = L_{\mathrm{LM}} + \alpha \cdot L_{\mathrm{balance}}
  \label{eq:aux-loss}
\end{equation}
其中辅助损失的形式为：
\begin{equation}
  L_{\mathrm{balance}} = N \cdot \sum_{i=1}^{N} f_i \cdot P_i
  \label{eq:balance-loss}
\end{equation}
这里 $f_i$ 是实际被分配到专家 $i$ 的 token 比例，
$P_i$ 是所有 token 对专家 $i$ 的平均路由概率。
$N$ 是专家数量，用于归一化。

直觉上，$f_i \cdot P_i$ 的乘积形式
鼓励两者同时尽可能均匀：
当所有专家均匀接收 token 时
（$f_i = P_i = 1/N$），
$\sum f_i P_i = N \times (1/N)^2 = 1/N$，达到最小值。
当负载高度不均衡时（所有 token 集中在一个专家），
$\sum f_i P_i = 1$，达到最大值。

Switch Transformer~\cite{fedus2022switch} 简化了这个公式，
使用 $\alpha = 0.01$ 作为默认系数，
并证明即使是 top-1 路由，
只要辅助损失设置得当，训练也能保持稳定。

但辅助损失有一个\textbf{根本矛盾}：
$\alpha$ 太大，负载均衡得到满足，
但辅助损失的梯度「污染」了语言建模的梯度——
路由器不再纯粹地为质量最优而路由，
而是在「质量」和「均衡」之间妥协，模型质量下降。
$\alpha$ 太小，辅助损失形同虚设，
专家坍缩照常发生。
在实践中，$\alpha$ 的调节需要大量实验，
且没有一个值能适用于训练的所有阶段。

\subsection{DeepSeek 的无辅助损失方案}

DeepSeek-V3~\cite{deepseekv3} 提出了一个优雅的替代方案：
\textbf{不加辅助损失}，
而是给每个专家的路由分数加一个\textbf{不参与梯度计算}的偏置项 $b_i$：
\begin{equation}
  s_i = W_{\mathrm{router}} \mathbf{x} + b_i
  \quad (\text{$b_i$ 不参与反向传播})
\end{equation}
偏置的更新规则极其简单：
如果第 $i$ 个专家过载，$b_i \leftarrow b_i - \gamma$；
如果空闲，$b_i \leftarrow b_i + \gamma$（$\gamma \approx 10^{-3}$）。

因为 $b_i$ 不在计算图中，
语言建模的梯度完全不受干扰——
训练目标纯粹是预测下一个 token，没有任何妥协。
2025 年芝加哥大学的理论论文证明了这种方法等价于
原始--对偶分配问题的在线求解，具有单调改进保证。

这个方法的优雅之处在于它将两个不同的优化问题\textbf{解耦}了：
语言建模（预测下一个 token）通过标准的梯度下降优化，
负载均衡通过偏置项的启发式规则调整。
两者互不干扰——语言建模的梯度不受负载均衡的影响，
负载均衡的调整也不需要通过梯度传播。
这就像一个工厂的生产线和排班表是分开管理的——
生产线的优化不需要考虑排班约束，
排班表的调整也不需要改变生产流程。

\begin{keybox}[负载均衡方法对比]
\begin{center}
\begin{tabular}{lccl}
  \hline
  \textbf{方法} & \textbf{需要辅助损失} & \textbf{影响 LM 梯度}
  & \textbf{代表模型} \\
  \hline
  辅助损失 & 是 & 是（$\alpha$ 敏感） & Switch Transformer \\
  Expert Choice & 否 & 部分 & Expert Choice Routing \\
  偏置项调节 & 否 & 否 & DeepSeek-V3 \\
  \hline
\end{tabular}
\end{center}
\end{keybox}

\section{DeepSeek 的三大创新}

DeepSeek-V3~\cite{deepseekv3} 在 MoE 上做了三项改变游戏规则的创新。

\subsection{细粒度专家}

传统 MoE 使用少量大专家（如 Mixtral 的 8 个专家，每个约 7B 参数）。
DeepSeek 使用 \textbf{256 个小专家}（每个仅约 0.5B 参数）
加 1 个\textbf{共享专家}（处理所有 token 的通用知识）：
\begin{equation}
  \mathbf{y} = \mathrm{SharedExpert}(\mathbf{x})
  + \sum_{i \in \mathrm{Top8}} g_i \cdot \mathrm{Expert}_i(\mathbf{x})
\end{equation}

256 个专家中选 8 个，组合数为
$\binom{256}{8} \approx 4.4 \times 10^{9}$——
相比 Mixtral 的 $\binom{8}{2} = 28$ 种组合，
组合空间扩大了 $10^{8}$ 倍。
这意味着模型可以为不同类型的输入激活完全不同的专家组合，
实现更精细的专业化。

为什么细粒度专家更好？三个原因：

\textbf{更灵活的组合}。
8 个大专家中选 2 个，
就像只有 8 门课的课程表中选 2 门——
选择空间很有限。
256 个小专家中选 8 个，
就像 256 门课中选 8 门——
可以精确匹配每个学生（token）的需求。

\textbf{更好的专业化}。
小专家被迫在一个更窄的领域内做到精通，
因为它们的参数量有限，无法成为「通才」。
大专家则可能在多个不相关的领域之间分散注意力。
实验表明，细粒度专家的专业化程度更高——
某些专家可能只处理 Python 代码中的列表推导，
另一些只处理数学证明中的符号推理。

\textbf{相似的计算成本}。
关键洞察：8 个小专家的总计算量
$\approx$ 2 个大专家的计算量（如果大小比例为 1:4）。
所以细粒度专家不增加每个 token 的计算成本——
它只是用更多的小模块替代了少数大模块。

\subsection{共享专家}

在上面的公式中，
$\mathrm{SharedExpert}(\mathbf{x})$ 是一个始终被激活的专家——
不管 token 是什么类型，它都参与计算。

为什么需要共享专家？
在纯路由的 MoE 中，
每个路由专家都需要独立学会一些「基础知识」
（如语法规则、常见词的语义等），
因为不能保证某个 token 每次都被路由到同一个专家。
这导致了跨专家的冗余——
256 个专家中可能有 200 个都各自学了一遍「the 后面通常跟名词」。

共享专家将这些通用知识集中到一个始终激活的模块中，
路由专家只需学习\textbf{差异化}的专业知识。
这既减少了冗余，又让路由专家能更专注于自己的领域。

从残差连接的视角看，
共享专家提供了一个强大的「默认输出」——
即使路由器把 token 送到了不太合适的专家，
共享专家的输出仍然保证了基本的质量。

\subsection{DeepEP：专用通信库}

MoE 的一个工程瓶颈是 All-to-All 通信：
每个 token 需要被发送到对应专家所在的 GPU，
计算完成后再发送回来。
这种细粒度的跨 GPU 通信模式不同于数据并行的 All-Reduce，
标准通信库的效率很低。

为什么 All-to-All 比 AllReduce 更难优化？
AllReduce 中每块 GPU 发送和接收的数据量是确定性的——
总是 $M/N$ 的数据。
All-to-All 中数据量取决于路由决策——
如果某个专家特别热门，
对应 GPU 可能收到 10 倍于平均水平的数据，
而其他 GPU 几乎收不到数据。
这种\textbf{不均匀的通信模式}
让带宽利用和延迟优化变得困难。

DeepSeek 开源的 DeepEP 提供两种模式：

\textbf{Normal 模式}用于训练和 prefill（高吞吐量优先）：
将大量 token 的路由请求打包成批次，
利用 NVLink（节点内）和 InfiniBand（节点间）的
分层通信策略——
同一节点内的专家通信走 NVLink（900\,GB/s），
跨节点的专家通信走 IB（400\,Gb/s），
两者可以并行进行。

\textbf{Low-Latency 模式}用于推理 decode
（通过 hook-based 通信--计算重叠实现零额外 SM 占用）。
在 decode 阶段，每个时间步只有 1 个 token 需要路由，
通信量很小但对延迟极其敏感。
DeepEP 利用 GPU 的硬件 hook 机制
将通信操作卸载到 NVLink/IB 控制器，
SM 完全不参与通信——
计算和通信真正同时进行，零开销。

\section{MoE 在实践中}

\subsection{Mixtral 8x7B：MoE 的开源先驱}

Mixtral 8x7B~\cite{jiang2024mixtral} 是 Mistral AI 发布的
第一个广泛使用的开源 MoE 模型。
它的规格说明了 MoE 的效率优势：
\begin{itemize}[noitemsep]
  \item \textbf{总参数}：47B（8 个 FFN 专家 $\times$ $\sim$5.5B + 非 FFN 参数）
  \item \textbf{活跃参数}：13B（每个 token 激活 2 个专家）
  \item \textbf{性能}：在大部分基准测试上
        匹配或超越 LLaMA 2 70B（稠密模型）
  \item \textbf{推理成本}：约为 LLaMA 2 70B 的 $1/5$
\end{itemize}

47B 总参数 vs 70B 稠密参数，
Mixtral 的参数总量反而更少——
但它用 MoE 的稀疏激活实现了类似甚至更好的性能，
同时推理时只需要 13B 的计算量。

\subsection{DeepSeek-V3：MoE 的巅峰}

DeepSeek-V3~\cite{deepseekv3} 将 MoE 推向了新的高度：
\begin{itemize}[noitemsep]
  \item \textbf{总参数}：671B
  \item \textbf{活跃参数}：37B（256 路由专家选 8 + 1 共享专家）
  \item \textbf{训练数据}：14.8T token
  \item \textbf{训练成本}：$\sim$\$5.6M（2048 块 H800，约 2 个月）
  \item \textbf{性能}：在多项基准测试上达到或超越 GPT-4o、Claude 3.5 Sonnet
\end{itemize}

$\$$5.6M 的训练成本令整个行业震惊——
它比同等性能的 LLaMA 3 405B 的估计训练成本低约 20 倍。
成本优势来自三方面叠加：
MoE 架构（37B 活跃 vs 405B 全激活）、
FP8 训练（减少 $\sim$50\% 内存和 $\sim$30\% 计算）、
以及精细的工程优化（DualPipe、DeepEP 等）。

\subsection{LLaMA 4：Meta 的 MoE 转型}
\label{sec:llama4-moe}

2025 年 4 月，Meta 发布了 LLaMA 4 系列——% NEW REF: llama4 = The Llama 4 Herd, Meta AI, 2025
这是 LLaMA 家族首次采用 MoE 架构，
标志着 Meta 从稠密模型全面转向 MoE。

LLaMA 4 包含三个变体：
\begin{center}
\begin{tabular}{lcccl}
  \hline
  \textbf{模型} & \textbf{总参数} & \textbf{活跃参数}
  & \textbf{专家数 / Top-K} & \textbf{特色} \\
  \hline
  Scout & 109B & 17B & 16 / 1 & 10M token 上下文 \\
  Maverick & 400B & 17B & 128 / 1 & 原生多模态 \\
  Behemoth & $\sim$2T & 288B & 16 / 2 & 训练中（未发布） \\
  \hline
\end{tabular}
\end{center}

LLaMA 4 的 MoE 设计与 DeepSeek 有显著差异：

\textbf{Top-1 路由}。
Scout 和 Maverick 都使用 Top-1 路由
（每个 token 只发给 1 个专家），
而不是 DeepSeek 的 Top-8。
这是一个大胆的选择——
Top-1 在 Switch Transformer 中曾被证明可行，
但需要更精确的路由器。
Meta 通过从 Behemoth（Top-2）的知识蒸馏
来弥补 Top-1 路由的信息损失。

\textbf{知识蒸馏}。
Scout 和 Maverick 都通过从 Behemoth 蒸馏获得了
显著的质量提升——这是一种``大模型训练小模型''的策略。
Behemoth 本身有 288B 活跃参数（Top-2 路由 16 个专家），
在 STEM 基准测试上超越了 GPT-4.5 和 Claude Sonnet 3.7。

\textbf{超长上下文}。
Scout 支持 \textbf{10M token 上下文窗口}——
目前最长的生产级上下文长度。
这通过高效的 iRoPE 位置编码
和分块注意力机制实现，
但也对 MoE 的路由一致性提出了更高要求。

\subsection{Qwen3-235B：阿里巴巴的 MoE 实践}

Qwen3-235B-A22B 是阿里巴巴 Qwen 团队% NEW REF: qwen3 = Qwen3 Technical Report, Qwen Team, arXiv 2025
发布的旗舰 MoE 模型，采用了与 DeepSeek 相似但不同的设计选择：

\begin{itemize}[noitemsep]
  \item \textbf{总参数}：235B，\textbf{活跃参数}：22B。
  \item \textbf{128 个路由专家，Top-8 路由}——
        与 DeepSeek-V3 的 256 专家 + Top-8 相比，
        专家数减半但每个专家更大。
  \item \textbf{94 层}，使用 GQA（64 头，4 KV 头）。
  \item 支持 128K token 上下文，通过 YaRN 位置编码外推。
\end{itemize}

Qwen3 的一个独特创新是\textbf{思考/非思考双模式}——
模型可以在``思考模式''（生成内部推理链）和
``直接回答模式''之间切换，
用户通过 prompt 控制推理预算。
这对 MoE 路由提出了新要求：
思考模式下的 token 可能需要不同的专家组合
来处理推理链中的中间步骤。

\subsection{Kimi K2：万亿参数的 MoE}

Moonshot AI 的 Kimi K2 将 MoE 推到了\textbf{万亿参数}量级：% NEW REF: kimik2_moe = Kimi K2: Open Agentic Intelligence, Kimi Team, arXiv 2025
\begin{itemize}[noitemsep]
  \item \textbf{总参数}：$\sim$1T（1.04T）。
  \item \textbf{活跃参数}：32B。
  \item \textbf{384 个路由专家，Top-8 路由}——
        比 DeepSeek-V3 更多的专家（384 vs 256），
        组合空间为 $\binom{384}{8} \approx 1.8 \times 10^{13}$。
  \item 使用 \textbf{MLA}（Multi-head Latent Attention）
        压缩 KV cache。
  \item 128K token 上下文。
\end{itemize}

Kimi K2 在 MoE 训练上的关键创新是
将 \textbf{Muon 优化器}引入 MoE 训练
（使用改进版 MuonClip），
并在 15.5T token 的全程训练中实现了``零 loss spike''。
这是一个重要的工程里程碑——
万亿参数 MoE 模型的训练稳定性
一直是业界最大的担忧之一。

\begin{keybox}[2025--2026 主要 MoE 模型对比]
\begin{center}
\begin{tabular}{lccccc}
  \hline
  \textbf{模型} & \textbf{总参数} & \textbf{活跃参数}
  & \textbf{专家数} & \textbf{Top-K} & \textbf{共享专家} \\
  \hline
  Mixtral 8x7B & 47B & 13B & 8 & 2 & 无 \\
  DeepSeek-V3 & 671B & 37B & 256 & 8 & 1 \\
  LLaMA 4 Maverick & 400B & 17B & 128 & 1 & 无$^\dagger$ \\
  LLaMA 4 Behemoth & $\sim$2T & 288B & 16 & 2 & 无$^\dagger$ \\
  Qwen3-235B & 235B & 22B & 128 & 8 & 未公开 \\
  Kimi K2 & 1.04T & 32B & 384 & 8 & 有 \\
  \hline
\end{tabular}
\end{center}

\small $^\dagger$\,LLaMA 4 未明确公开共享专家设计。
\end{keybox}

从表中可以看到一个清晰的趋势：
\textbf{细粒度小专家 + 高 Top-K} 成为主流设计
（DeepSeek、Qwen3、Kimi K2 都使用 Top-8 + 100 个以上的专家）。
LLaMA 4 的 Top-1 路由 + 知识蒸馏是一条独特的路线，
但 Behemoth（未发布的旗舰模型）也回到了 Top-2。

\section{MoE 推理的工程挑战}
\label{sec:moe-inference}

MoE 的训练效率优势在推理端面临严峻的工程挑战。

\subsection{内存瓶颈与专家卸载}

MoE 模型的推理面临一个核心矛盾：
\textbf{每次只激活 5\%--10\% 的参数，
但 100\% 的参数都必须驻留在可访问的存储中}。
DeepSeek-V3 的 671B 参数在 BF16 下需要 $\sim$1.3\,TB 内存，
至少需要 17 块 H100（80\,GB）才能装下——
但每个 token 实际只需要 37B 参数（$\sim$74\,GB）的计算。

\textbf{专家卸载}（Expert Offloading）将不活跃的专家
从 GPU 显存卸载到 CPU 内存甚至 SSD，
只在需要时加载到 GPU：

\textbf{FineMoE}~(2026) 提出了\textbf{细粒度}的卸载策略：% NEW REF: finemoe2026 = FineMoE: Taming Latency-Memory Trade-Off in MoE-Based LLM Serving via Fine-Grained Expert Offloading, Yu et al., EuroSys 2026
不是以整个专家为单位进行卸载/加载
（粒度太粗——一个专家可能有几百 MB），
而是将每个专家进一步拆分为更小的``子专家''块，
按需加载。
关键洞察是：专家的使用频率遵循\textbf{长尾分布}——
少数``热门专家''被频繁调用，大部分专家很少使用。
FineMoE 将热门专家常驻 GPU，
冷门专家按子块粒度从 CPU 按需传输。

\textbf{ADEPT}（2026）进一步引入了% NEW REF: adept2026 = Two-Stage Expert Offloading for Domain-Aware MoE Inference, Kim et al., IEEE Access 2026
\textbf{领域感知的预取}策略：
根据输入文本的领域特征（代码、数学、自然语言等）
预判哪些专家可能被激活，提前将它们从 CPU 加载到 GPU。
这将专家加载操作从推理的关键路径上移除，
有效隐藏了 PCIe 总线的 I/O 延迟。

\subsection{MoE 模型的量化}

量化是降低 MoE 内存占用的另一条路径——
但 MoE 的量化比稠密模型更具挑战性。

\textbf{DynaExq}~(2025) 提出了% NEW REF: dynaexq2025 = Dynamic Expert Quantization for Scalable MoE Inference, Chu et al., arXiv 2025
\textbf{动态混合精度}的量化方案：
不同专家使用不同的量化位宽。
热门专家（频繁使用）保持较高精度（如 8 bit），
冷门专家使用更激进的压缩（如 4 bit 甚至 2 bit），
因为它们对最终输出的贡献权重较小。
DynaExq 在运行时根据路由统计量
动态调整专家的量化位宽，
在精度和内存之间找到最优平衡。

\textbf{MiLo}（Mixture of Low-Rank Compensators）% NEW REF: milo2025 = MiLo: Efficient Quantized MoE Inference with Mixture of Low-Rank Compensators, Huang et al., arXiv 2025
则在量化后为每个专家附加一个\textbf{低秩补偿器}——
用少量的额外参数修正量化带来的误差。
对于 4-bit 量化的 MoE 模型，
MiLo 可以恢复 90\%+ 的原始精度，
而额外的参数开销不到原始专家的 5\%。

\subsection{批处理的矛盾}

MoE 推理还面临一个``批处理悖论''：
\textbf{批处理是推理吞吐量的关键，
但它抵消了 MoE 的稀疏性优势}。

在单请求推理中，每个 token 只激活 $K$ 个专家，
MoE 的计算节省是真实的。
但在批处理（batched inference）中，
不同请求的 token 可能路由到不同的专家——
当 batch size 足够大时，
所有 $N$ 个专家都会被至少一个 token 激活，
等价于运行一个\textbf{完整的稠密模型}。
更糟的是，MoE 的 All-to-All 通信开销在稠密模型中不存在。

\textbf{Lynx}~(2025) 提出了一种``专家重映射''方案：% NEW REF: lynx2025 = Lynx: Enabling Efficient MoE Inference via Dynamic Batch-Aware Expert Selection, arXiv 2025
在运行时动态地将 batch 中不同请求的 token
重新分配到重叠度更高的专家子集上，
减少实际需要激活的专家总数。
这在保持 $<$1\% 精度损失的前提下，
将推理吞吐量提升了 $1.23\times$。

\textbf{MoE-SpeQ}~(2025) 则将推测解码% NEW REF: moespeq2025 = MoE-SpeQ: Speculative Quantized Decoding with Proactive Expert Prefetching and Offloading for MoE, Wang et al., arXiv 2025
与专家预取结合：
在主模型验证候选 token 的同时，
利用路由器的预判结果预加载下一步可能需要的专家。
推测解码本身减少了生成步数，
专家预取则隐藏了每步的加载延迟——
两种优化相互增强。

\subsection{MoE 的优势与挑战}

\begin{corebox}[MoE vs Dense：何时选择哪种架构]
\textbf{MoE 的优势}：
\begin{itemize}[noitemsep]
  \item 知识容量与推理成本解耦——大容量、低推理成本
  \item 对于广泛部署的模型，推理经济性是压倒性优势
  \item 专家专业化带来的任务适配能力
\end{itemize}

\textbf{MoE 的挑战}：
\begin{itemize}[noitemsep]
  \item \textbf{内存占用}：所有专家必须加载到内存中，
        即使每次只激活一小部分。
        671B MoE 模型的权重（BF16）约 1.3\,TB——
        需要至少 17 块 H100（80\,GB）才能放下。
  \item \textbf{微调困难}：
        LoRA 等参数高效微调方法
        在 MoE 上的应用并不直接——
        应该给每个专家加 LoRA 适配器，
        还是只给共享部分加？
        给所有 256 个专家都加 LoRA 会导致参数量爆炸，
        只给部分专家加则可能破坏路由的平衡。
  \item \textbf{训练复杂度}：
        负载均衡、通信优化、容量因子调节——
        MoE 的超参数空间比稠密模型大得多。
  \item \textbf{推理延迟}：
        虽然 FLOP 更少，但 All-to-All 通信引入了额外延迟。
        在单 GPU 推理中，MoE 的延迟优势可能被
        更多的内存访问抵消
        （所有专家的参数都在内存中，cache 命中率更低）。
\end{itemize}

\textbf{经验法则}：
当目标模型的知识容量需要 $>$100B 参数，
且推理成本是重要考量时，选择 MoE。
当模型规模在 7B--14B，
且部署环境受限（单 GPU、边缘设备）时，选择 Dense。
\end{corebox}

\begin{whybox}[为什么 MoE 是 2026 年的标准选择？]
答案是\textbf{推理经济学}。训练一次，推理无数次。
MoE 使得你可以用 37B 的推理成本获得 671B 的知识容量。
当你服务数百万用户、每天处理数十亿 token 时，
这个 $\sim$20$\times$ 的推理计算节省是巨大的。
这就是为什么 2026 年所有 frontier 开源模型
（DeepSeek-V3/V4、LLaMA 4、Qwen3-235B、Mistral Large 3）
都采用了 MoE 架构。

值得注意的是，MoE 并不总是最优选择。
对于参数量在 7B--14B 的小模型，
Dense（稠密）架构仍然占主导——
因为小模型的参数本来就不多，
再拆分成多个专家后每个专家太小，
专业化程度不够。
而且小模型通常部署在单块 GPU 甚至 CPU 上，
MoE 的 All-to-All 通信开销在单设备上没有意义。

MoE 的真正优势在大规模模型上：
当你需要 $>$100B 的知识容量，
但推理成本必须控制在 $\sim$20B 参数以内时，
MoE 是唯一的选择。
\end{whybox}

\section{2025--2026 年 MoE 架构的爆发}
\label{sec:moe-explosion}

2025 年下半年至 2026 年初，
MoE 从少数模型的架构选择演变为前沿中国模型的\textbf{标准范式}。
短短几个月内，至少 9 个主要模型采用了 MoE 架构，
总参数规模从 80B 跨越到 1T，
而活跃参数却压缩到惊人的低水平。

\begin{center}
\begin{tabular}{llrrll}
  \hline
  \textbf{模型} & \textbf{发布日期} & \textbf{总参数} & \textbf{活跃参数}
  & \textbf{专家配置} & \textbf{亮点} \\
  \hline
  DeepSeek V4$^\star$ & 2026 年 4 月（预期） & $\sim$1T & $\sim$320B
  & -- & + Engram + mHC \\
  Kimi K2.5 & 2026 年 1 月 & 1T & 32B
  & 384（8 active） & Agent Swarm \\
  GLM-5 & 2026 年 2 月 & 744B & 40B
  & -- & 零 NVIDIA 训练 \\
  Qwen 3.5 & 2026 年 2 月 & 397B & 17B
  & -- & Gated Delta Networks \\
  Qwen3-Next & 2025 年 9 月 & 80B & 3B
  & 512（10+1 active） & 超稀疏激活 \\
  Hunyuan 2.0 & 2025 年 12 月 & 406B & 32B
  & -- & 腾讯 \\
  MiniMax M2.5 & 2026 年 2 月 & 230B & 10B
  & -- & SWE-bench 80.2\% \\
  Step 3.5 Flash & 2026 年 2 月 & 196B & 11B
  & -- & 单卡 RTX 4090 可运行 \\
  豆包 Seed 2.0 & 2026 年 2 月 & -- & --
  & -- & AIME 98.3\% \\
  \hline
\end{tabular}

\small $^\star$\,DeepSeek V4 截至 2026 年 3 月尚未正式发布，
参数为预期值。
\end{center}

这张表格揭示了 2025--2026 年 MoE 生态的三个核心趋势。

\subsection{超稀疏激活}

MoE 的稀疏度正在被推向新的极端。
DeepSeek-V3 的激活率约为 5.5\%（37B / 671B），
已经令人印象深刻——
但 \textbf{Qwen3-Next} 将这个数字压低到了
\textbf{3.75\%}（3B / 80B）。

Qwen3-Next 使用了 512 个路由专家，
每个 token 仅激活 10+1 个（10 个路由 + 1 个共享）。
80B 的模型只需要 3B 的推理计算量——
这意味着一个具有 80B 参数知识容量的模型，
推理成本仅相当于一个小型 3B 稠密模型。
激活比例从 DeepSeek-V3 的 1:18 演进到 Qwen3-Next 的 1:27，
超稀疏激活正在成为 MoE 设计的新前沿。

\subsection{消费级硬件部署}

\textbf{Step 3.5 Flash} 证明了 196B 参数的 MoE 模型
可以在\textbf{单块 RTX 4090}（24\,GB 显存）上运行。
11B 的活跃参数经过 4-bit 量化后仅需约 6\,GB，
其余参数通过专家卸载动态加载。

这打破了``大模型需要大集群''的思维定式——
MoE 的稀疏激活特性使得总参数量不再是部署的瓶颈，
活跃参数量才是。
当活跃参数降到 10--17B 范围时，
消费级 GPU 就足以提供可接受的推理体验。

\subsection{总参数与活跃参数的分离}

从上表可以清晰看到，
行业正在收敛到一个模式：
\textbf{10--40B 的活跃参数 + 200B--1T 的总参数}。
这意味着知识容量（总参数）和推理成本（活跃参数）
被彻底解耦——
模型可以拥有巨大的知识储备，
但每次推理只调用极小一部分。

DeepSeek V4（预期）的 $\sim$320B 活跃参数是一个例外——
它选择了更高的活跃参数来支持原生多模态理解与生成、
Engram 长期记忆和 mHC（流形约束超连接）等创新架构，
代价是更高的推理成本。
这也表明：活跃参数量的选择
最终取决于模型需要\textbf{同时在线}的能力复杂度。

\section{MoE 理论的新进展}
\label{sec:moe-theory-2026}

伴随 MoE 在实践中的爆发，
理论层面也涌现了重要的新成果。

\textbf{Optimal Expert-Attention Allocation}% NEW REF: expertattention2026 = Optimal Expert-Attention Allocation in MoE, arXiv:2603.10379, 2026
（arXiv:2603.10379，2026 年 3 月）
将神经 Scaling Laws 扩展到了 MoE 架构。
传统 Scaling Laws 将模型视为同质的参数集合，
而 MoE 模型中存在两种本质不同的计算单元：
\textbf{专家层}（负责知识检索和领域特化）和
\textbf{注意力层}（负责上下文交互和信息整合）。
这篇论文研究了在固定计算预算下，
如何最优地分配计算资源到专家层和注意力层——
结论是最优分配比例与模型规模和任务类型密切相关，
且与直觉中``注意力更重要''的假设不完全一致。

\textbf{LatentMoE}（arXiv:2601.18089，2026 年 1 月）% NEW REF: latentmoe2026 = LatentMoE: Toward Optimal Accuracy per FLOP and Parameter in MoE, arXiv:2601.18089, 2026
探索了在 MoE 中实现\textbf{每 FLOP 和每参数最优精度}的路径。
它通过在潜在空间（latent space）中执行路由和专家计算，
减少了路由决策的维度，
同时在理论上保证了与标准 MoE 相同的表达能力。
这为进一步压缩 MoE 的推理成本提供了理论基础。

\textbf{NVIDIA Megatron Core MoE}（arXiv:2603.07685，2026 年 3 月）% NEW REF: megatron_core_moe_2026
则从系统层面解决了 MoE 训练中
\textbf{内存、通信与计算的三方耦合}问题。
在大规模 MoE 训练中，
增大专家数量会同时增加内存占用（专家权重）、
通信量（All-to-All 传输）和计算模式的碎片化——
这三者的交互使得单独优化任一维度可能恶化其他维度。
Megatron Core MoE 通过全栈协同设计，
将专家并行、张量并行和流水线并行
纳入统一的抽象层，系统化地平衡这些约束。


%% ============================================================
% NEW REF: jacobs1991moe = Adaptive Mixtures of Local Experts, Jacobs et al., Neural Computation 1991
% NEW REF: lepikhin2021gshard = GShard: Scaling Giant Models with Conditional Computation, Lepikhin et al., ICLR 2021
% NEW REF: zhou2022expertchoice = Mixture-of-Experts with Expert Choice Routing, Zhou et al., NeurIPS 2022
%% --- NEW REFERENCES (March 2026 update) ---
% NEW REF: dynamoe2026 = DynaMoE: Dynamic Token-Level Expert Activation with Layer-Wise Adaptive Capacity, Guelmez, arXiv 2026
% NEW REF: dirmoe2026 = DirMoE: Differentiable Probabilistic Router for Mixture-of-Experts, ICLR 2026
% NEW REF: l2r2026 = L2R: Low-Rank and Lipschitz-Controlled Routing for Mixture-of-Experts, Yang et al., arXiv 2026
% NEW REF: llama4 = The Llama 4 Herd, Meta AI, 2025
% NEW REF: qwen3 = Qwen3 Technical Report, Qwen Team, arXiv 2025
% NEW REF: kimik2_moe = Kimi K2: Open Agentic Intelligence, Kimi Team, arXiv 2025
% NEW REF: finemoe2026 = FineMoE: Taming Latency-Memory Trade-Off in MoE-Based LLM Serving, Yu et al., EuroSys 2026
% NEW REF: adept2026 = Two-Stage Expert Offloading for Domain-Aware MoE Inference, Kim et al., IEEE Access 2026
% NEW REF: dynaexq2025 = Dynamic Expert Quantization for Scalable MoE Inference, Chu et al., arXiv 2025
% NEW REF: milo2025 = MiLo: Efficient Quantized MoE Inference with Low-Rank Compensators, Huang et al., arXiv 2025
% NEW REF: lynx2025 = Lynx: Enabling Efficient MoE Inference via Dynamic Batch-Aware Expert Selection, arXiv 2025
% NEW REF: moespeq2025 = MoE-SpeQ: Speculative Quantized Decoding with Expert Prefetching for MoE, Wang et al., arXiv 2025
% NEW REF: expertattention2026 = Optimal Expert-Attention Allocation in MoE, arXiv:2603.10379, 2026
% NEW REF: latentmoe2026 = LatentMoE: Toward Optimal Accuracy per FLOP and Parameter in MoE, arXiv:2601.18089, 2026
% NEW REF: megatron_core_moe_2026 = NVIDIA Megatron Core MoE: Full-Stack System Co-Design for MoE Training, NVIDIA, arXiv:2603.07685, 2026
